\documentclass[11pt]{article}
\usepackage[margin=1in]{geometry}
\usepackage{amsmath,amssymb,amsthm,enumitem,graphicx,xcolor,hyperref,microtype}
\usepackage[T1]{fontenc}
\usepackage{lmodern}
\hypersetup{colorlinks=true,linkcolor=blue!55!black,urlcolor=blue!55!black,citecolor=blue!55!black}
\newtheorem{theorem}{Theorem}
\newtheorem{lemma}{Lemma}
\newtheorem{corollary}{Corollary}
\newcommand{\C}{\mathbb C}
\newcommand{\R}{\mathbb R}
\newcommand{\N}{\mathbb N}
\newcommand{\cnull}{c_0}
\title{No Lipschitz Tube-Holomorphic Separator\\on a Banach Space Containing $c_0$}
\author{Automated mathematics run \texttt{fa\_banach\_001}, lane 17}
\date{11 August 2026}

\begin{document}
\maketitle

\begin{center}
\fbox{\parbox{0.91\linewidth}{\textbf{Status.}
Candidate full solution, likely valid, requiring human review.  The theorem
answers negatively the explicit $c_0$ version of the Lipschitz-extension
hypothesis in source Theorem 2 and strengthens it to every space containing
an isomorphic copy of $c_0$.}}
\end{center}

\begin{abstract}
Azagra, Fry, and Keener ask whether $c_0$ has a Lipschitz real-analytic
separating function of the kind used in their approximation theorem: its
holomorphic extension is Lipschitz on a complex tube of fixed positive width.
We prove that it does not.  More generally, no real Banach space containing an
isomorphic copy of $c_0$ admits such a separator.  Continuous homogeneous
polynomials have summable diagonal values on a $c_0$-basic sequence.  The
Taylor expansion therefore makes the holomorphic extension vanish
asymptotically along the corresponding coordinate lines near zero.  Its
global Lipschitz bound makes those line restrictions a normal family on one
common strip, so the identity theorem propagates the convergence to every
fixed real point.  This contradicts linear separation at infinity.
\end{abstract}

\section{Source question}

The source is D. Azagra, R. Fry, and L. Keener, \emph{Real analytic
approximation of Lipschitz functions on Hilbert space and other Banach
spaces}, arXiv:1005.1050v5, published in J. Funct. Anal. \textbf{262} (2012),
124--166 \cite{AFK}.  It calls $Q:X\to[0,\infty)$ a separating function if
$Q(0)=0$ and there are $m,M>0$ such that
\[
 Q(x)\ge m\|x\|\qquad(\|x\|\ge M).
\]
Theorem 2 assumes that $Q$ has a \emph{Lipschitz holomorphic extension}
$\widetilde Q$ on a fixed-width tube
\[
 \Omega_\delta=\{x+iy\in X_{\C}:x,y\in X,\ \|y\|<\delta\}.
\]
Immediately afterward, on PDF page 4, the authors write:
\begin{quote}
``We do not know if the converse is true. A natural related question is: does
the space $c_0$ have such a Lipschitz separating function?''
\end{quote}

\begin{figure}[ht]
\centering
\includegraphics[width=0.94\linewidth]{figures/open_problem_crop.png}
\caption{Theorem 2 and the $c_0$ question in arXiv:1005.1050v5, PDF page 4.}
\end{figure}

We answer the explicit Theorem 2 version: the extension itself is globally
Lipschitz on $\Omega_\delta$.

\section{Proof intuition}

Let $(u_n)$ be a sequence equivalent to the canonical basis of $c_0$.
The key algebraic fact is that for every continuous homogeneous polynomial
$P$, the diagonal sequence $(P(u_n))$ is summable.  One sees this by inserting
independent roots of unity into finite sums $\sum z_nu_n$: averaging kills all
mixed monomials and retains only $P(u_n)$.

Now let $F$ be the hypothetical holomorphic extension.  Its Taylor expansion
at zero and diagonal summability imply $F(tu_n)\to0$ for small real $t$.
The scalar functions $f_n(z)=F(zu_n)$ are holomorphic on one common horizontal
strip.  Because $F$ is Lipschitz on the whole tube, they are uniformly bounded
on compact subsets of that strip.  Montel's theorem gives locally uniform
subsequential limits; every limit vanishes on the small real interval and is
therefore identically zero.  Thus $F(Tu_n)\to0$ for every fixed real $T$.
Choosing $T$ large contradicts the separator inequality, uniformly in $n$.

\section{The diagonal polynomial lemma}

We use the following form, which makes the later strengthening automatic.

\begin{lemma}[summable diagonals]\label{lem:diagonal}
Let $E$ be a complex Banach space and let $(u_n)\subset E$ satisfy
\[
 \left\|\sum_{n=1}^N z_nu_n\right\|
 \le B\max_{1\le n\le N}|z_n|
 \tag{1}\label{eq:upper}
\]
for all $N$ and all complex scalars $z_1,\ldots,z_N$.  If $P:E\to\C$ is a
continuous $k$-homogeneous polynomial, then
\[
 \sum_{n=1}^\infty |P(u_n)|<\infty.
\]
For $k\ge2$, in fact
$\sum_n|P(u_n)|\le B^k\|P\|$, where
$\|P\|=\sup_{\|x\|\le1}|P(x)|$.
\end{lemma}

\begin{proof}
For $k=1$, $P$ is linear.  Given $N$, choose unit scalars $\lambda_n$ so that
$\lambda_nP(u_n)=|P(u_n)|$.  Equation \eqref{eq:upper} gives
\[
 \sum_{n=1}^N|P(u_n)|
 =P\left(\sum_{n=1}^N\lambda_nu_n\right)
 \le B\|P\|.
\]

Let $k\ge2$ and let $A$ be the symmetric $k$-linear form associated with
$P$.  Choose unit scalars $\lambda_n$ such that
$\lambda_n^kP(u_n)=|P(u_n)|$.  Let $\xi_1,\ldots,\xi_N$ be independent,
uniformly distributed $k$th roots of unity.  Expanding
$P(\sum_n\lambda_n\xi_nu_n)$ in terms of $A$, a monomial survives expectation
only if the multiplicity of every $\xi_n$ is divisible by $k$.  Since the
total degree is $k$, the surviving monomials are precisely those whose $k$
indices are all equal.  Hence
\[
 \sum_{n=1}^N|P(u_n)|
 =\mathbb E\,P\left(\sum_{n=1}^N\lambda_n\xi_nu_n\right).
\]
Taking absolute values and using \eqref{eq:upper} yields the claimed bound,
uniformly in $N$.
\end{proof}

\section{Nonexistence theorem}

\begin{theorem}\label{thm:main}
Let $X$ be a real Banach space containing an isomorphic copy of $c_0$, and
let $X_{\C}$ be a complexification.  There do not exist
$\delta,L,m,M>0$ and a holomorphic function
\[
 F:\Omega_\delta=\{x+iy\in X_{\C}:\|y\|<\delta\}\longrightarrow\C
\]
such that:
\begin{enumerate}[label=\textup{(\roman*)}]
\item $F$ is $L$-Lipschitz on $\Omega_\delta$;
\item $Q=F|_X$ is real-valued, $Q(0)=0$, and $Q\ge0$;
\item $Q(x)\ge m\|x\|$ whenever $\|x\|\ge M$.
\end{enumerate}
\end{theorem}

\begin{proof}
Choose an isomorphic embedding $T:c_0\to X$ and put $u_n=Te_n$.  After
complexification there are constants $a,b>0$ such that
\[
 \|u_n\|\ge a,
 \qquad
 \left\|\sum_{n=1}^N z_nu_n\right\|
 \le b\max_n|z_n|
 \tag{2}\label{eq:c0}
\]
for real vectors in the first inequality and complex finite scalar families
in the second.

Write the Taylor expansion of $F$ at zero as
\[
 F(w)=\sum_{k=0}^\infty P_k(w),
\]
where $P_k$ is a continuous complex $k$-homogeneous polynomial and
$P_0=F(0)=0$.  Since $F$ is holomorphic and continuous at zero, there is a
ball about zero on which it is bounded, and its Taylor series converges
uniformly on every sufficiently smaller ball.  Choose $\rho>0$ so small that
the series converges uniformly whenever $\|w\|\le2b\rho$.

Fix a real $t$ with $|t|\le\rho$.  Lemma \ref{lem:diagonal}, applied to
\eqref{eq:c0}, gives $P_k(u_n)\to0$ for every fixed $k$.  Thus every fixed
finite Taylor sum at $tu_n$ tends to zero.  The uniform Taylor tail on
$\|tu_n\|\le b\rho$ is independent of $n$, so
\[
 F(tu_n)\longrightarrow0
 \qquad (|t|\le\rho).
 \tag{3}\label{eq:localzero}
\]

For every $n$, define $f_n(z)=F(zu_n)$.  Since $\|u_n\|\le b$, all $f_n$ are
holomorphic on the common strip
\[
 S=\{z\in\C:|\operatorname{Im}z|<\delta/b\}.
\]
The Lipschitz hypothesis and $F(0)=0$ give, for $z\in S$,
\[
 |f_n(z)|\le L\|zu_n\|\le Lb|z|.
 \tag{4}\label{eq:normal}
\]
Therefore $(f_n)$ is a normal family on $S$.  Take any subsequence.  Montel's
theorem provides a further subsequence converging locally uniformly to a
holomorphic function $f$ on $S$.  By \eqref{eq:localzero}, $f$ vanishes on
$[-\rho,\rho]$, and the identity theorem gives $f\equiv0$ on $S$.

It follows that the full sequence satisfies $f_n(z)\to0$ at every fixed
$z\in S$: otherwise a subsequence bounded away from zero at some $z$ would
have a locally uniformly convergent subsubsequence, whose limit we just proved
must be zero.

Finally choose a fixed real $R\ge M/a$.  Then
$\|Ru_n\|\ge aR\ge M$, and the separating inequality gives
\[
 f_n(R)=Q(Ru_n)\ge m\|Ru_n\|\ge maR>0
\]
for every $n$.  This contradicts $f_n(R)\to0$.
\end{proof}

\begin{corollary}
The space $c_0$ has no separating function whose holomorphic extension is
globally Lipschitz on a complex tube of fixed positive width.  Thus the answer
to the explicit Theorem 2 version of the source question is negative.
\end{corollary}

\section{Scope, comparison, and novelty status}

\paragraph{Why the classical $c_0$ separator does not conflict.}
The function
\[
 d(x)=\sum_{n=1}^\infty x_n^{2n}
\]
is a standard uniformly analytic separating function on $c_0$
\cite{BoisoHajek,Mytrofanov}.  It is not globally Lipschitz: on the coordinate
line $te_n$, its derivative has size $2n|t|^{2n-1}$.  The theorem concerns the
strictly stronger Lipschitz-tube hypothesis needed in source Theorem 2.

\paragraph{Interpretive limitation.}
The source's Theorem 2 explicitly says that the \emph{holomorphic extension}
is Lipschitz on the tube, and this is the property proved impossible here.
If the phrase ``such a Lipschitz separating function'' were instead intended
to require only that $Q$ be Lipschitz on the real space while allowing its
tube extension to have arbitrary growth, estimate \eqref{eq:normal} would no
longer follow and that weaker variant remains outside this packet.

\paragraph{Bounded novelty audit.}
The run indexes contained no duplicate.  OpenAlex listed thirteen citing
works for the published source as of 11 August 2026.  Their titles and the
available relevant full texts were checked.  Johanis's simpler proof
\cite{Johanis} retains the separating-polynomial hypothesis.  Mytrofanov's
survey \cite{Mytrofanov} and a 2021 symmetric-approximation paper discuss the
weaker uniformly analytic $c_0$ separator above, not a globally Lipschitz tube
extension.  Exact-phrase and close-variant searches found no direct answer.
Novelty confidence is moderate-to-high under this bounded audit, not an
exhaustive priority claim.

\paragraph{Human-review recommendation.}
Prioritize review.  Check the root-of-unity extraction in Lemma
\ref{lem:diagonal}, the normal-family step, and the intended interpretation of
the source's word ``such.''  No conditional mathematical dependency remains
for the explicit Theorem 2 version.

\begin{thebibliography}{9}
\bibitem{AFK}
D. Azagra, R. Fry, and L. Keener,
\emph{Real analytic approximation of Lipschitz functions on Hilbert space and
other Banach spaces}, J. Funct. Anal. \textbf{262} (2012), 124--166;
arXiv:1005.1050v5.

\bibitem{BoisoHajek}
M. Cepedello Boiso and P. H\'ajek,
\emph{Analytic approximations of uniformly continuous functions in real
Banach spaces}, J. Math. Anal. Appl. \textbf{256} (2001), 80--98.

\bibitem{Johanis}
M. Johanis,
\emph{A simple proof of the approximation by real analytic Lipschitz
functions}, J. Math. Anal. Appl. \textbf{388} (2012), 1--7.

\bibitem{Mytrofanov}
M. A. Mytrofanov,
\emph{Separating polynomials, uniform analytical and separating functions},
Carpathian Math. Publ. \textbf{7} (2015), 197--208.
\end{thebibliography}

\end{document}
